% Derived only from the existing su2_audit.json and su2_benchmark.json.
% No experiments were rerun for this subsection.

\subsection{SU(2) compilation and decision experiments}
\label{subsec:su2-experiments}

The recorded experiments test the exact formula construction and illustrate
the practical degree--variable tradeoff with the exact NLSAT method\cite{JovanovicDeMoura2012,ProgrammingZ3}.  They do not measure the running
time of the theoretical real-decision algorithm.  The complete records are
\path{results/su2_audit.json} and
\path{results/su2_benchmark.json}; the corresponding scripts are
\path{su2_research/audit_su2.py} and
\path{su2_research/benchmark_su2.py}.

\paragraph{Finite correctness audit.}
With seed $20261008$, the audit sampled 48 accepted one-component braid
closures on two to four strands, using words of one to seven letters.
Every third accepted diagram was mirrored.  An independently implemented
reduced $\mathbb F_2$ Khovanov cube classified 39 samples as unknots
and nine as knotted.  Bridge compilation completed 42 of 48 queries:
36 unknot decisions and six knottedness decisions.  Its six timeouts
comprised three samples of each oracle class.  Defining-generator Tietze
elimination followed by compilation completed all 48 queries.  Thus 90 of
96 queries were definitive, with no disagreement among those 90.
The audit identifies its exact backend as Z3 5.1.0.

These are 48 sampled entries, containing 38 distinct stored PD arrays,
and do not constitute 48 distinct knot types or an independent test-set
estimate of general success probability.  The saved PD arrays specify
the actual inputs: the stored source braid words precede the optional
mirroring step.  The seed and script reproduce that step.
For the 36 sampled entries with at most five crossings, comprising 26
distinct PD arrays, exhaustive counts of homomorphisms to $S_3$ agreed
for the Wirtinger, bridge, and simplified presentations.  This checks
concrete relation-order and substitution behavior; agreement for one
finite target group does not establish presentation equivalence in
general.  The Tietze and bridge proofs supply that justification.

The same audit checked the prime-pretzel family at
$m=3,5,7,11,21,51$, hence at $9,15,21,33,63,153$ crossings.
The computed Fox $3$-coloring dimensions were respectively
$3,5,7,11,21,51$, as predicted by
Proposition~\ref{prop:su2-prime-rank-obstruction}.  These finite checks
support the construction and indexing; the proposition's proof supplies
the statement for all odd $m$.

\paragraph{Presentation choice and peak degree.}
The benchmark makes three repetitions per presentation on seven fixed
examples, with a requested solver allowance of $250$ milliseconds.
Table~\ref{tab:su2-recorded-bench} reports representative cases.
Here $k$ is the actual number of real variables after the safe gauge and
meridional trace sharing, and $d_{\rm eq}$ is the maximum degree of an
individual compiled equality before squaring.  The implementation sends
the separate equations to NLSAT in these measurements.  Its sum-of-squares
aggregate is available separately for the constant-atom theoretical bound.

\begin{table}[htbp]
\centering
\small
\begin{tabular}{llrrrrl}
\hline
Input & Presentation & $k$ & $d_{\rm eq}$ & Compile ms & Query ms & Result\\
\hline
Figure-eight & Wirtinger & 10 & 4 & 0.702 & 273.510 & I\\
 & Bridge & 10 & 2 & 0.416 & 270.820 & I\\
 & Tietze & 4 & 10 & 1.716 & 26.103 & K\\
$T(2,9)$ & Wirtinger & 25 & 4 & 1.866 & 283.346 & I\\
 & Bridge & 25 & 2 & 0.972 & 285.034 & I\\
 & Tietze & 4 & 18 & 11.411 & 72.254 & K\\
Coxeter unknot, $n=9$ & Wirtinger & 25 & 4 & 1.519 & 280.035 & I\\
 & Bridge & 2 & 2 & 0.323 & 0.983 & U\\
 & Tietze & 2 & 2 & 0.319 & 1.232 & U\\
\hline
\end{tabular}
\caption{Recorded medians of three repetitions, in milliseconds.  K:
nonabelian representation; U: no nonabelian representation; I:
inconclusive timeout.  All three repetitions in each row had the displayed
status.  The Coxeter example is the closure of
$\sigma_1\cdots\sigma_9\in B_{10}$.  Timeout durations are reported
as censored runs and are not completion times.}
\label{tab:su2-recorded-bench}
\end{table}

On $T(2,9)$, Tietze elimination reduced $k$ from 25 to four while
raising $d_{\rm eq}$ to 18, with aggregate degree bound 36.
Its compilation took longer, but its query completed within the requested
allowance while the other two formulations did not.  This demonstrates
why generator elimination must be assessed together with degree and
support growth.  It gives no finite speedup ratio against a censored run.
Across all 63 benchmark queries, 39 were definitive and 24 inconclusive.

\paragraph{Compressed powers.}
A separate construction-only experiment uses the general presentation
$\langle x_1,x_2\mid x_1^{2^b}\rangle$, represented by $b$ squaring
gates.  It has no knot-group provenance and was not assigned a knot
verdict.  Table~\ref{tab:su2-recorded-powers} shows that the compiler
retained the circuit, producing constant-degree equations without
expanding $2^b$ letters.  The number of stored terms sums the supports
of the individual equations and the nonabelian polynomial; it is not the
support of a fully collected sum-of-squares aggregate.

\begin{table}[htbp]
\centering
\small
\begin{tabular}{rrrrr}
\hline
$b$ & Real variables & Stored terms & $d_{\rm eq}$ & Compile ms\\
\hline
8 & 37 & 97 & 2 & 0.260\\
32 & 133 & 361 & 2 & 0.815\\
128 & 517 & 1417 & 2 & 3.160\\
512 & 2053 & 5641 & 2 & 13.622\\
2048 & 8197 & 22537 & 2 & 54.236\\
\hline
\end{tabular}
\caption{Three-run compiler medians for an already constructed word
circuit.  Expanded word length is $2^b$; aggregate degree bound is four
throughout.  These measurements exclude WordArena construction and make
no claim about solver performance on the resulting formulas.}
\label{tab:su2-recorded-powers}
\end{table}

\paragraph{Interpretation and trust boundary.}
Both scripts enable a trace-zero witness probe before the unrestricted
query. A satisfying assignment to this restricted query is valid for
the original formula; an unsatisfiable or unknown probe is discarded by
this implementation. The known completeness theorem for verified knot
meridians, Theorem~\ref{thm:su2-traceless}, would permit a stronger
knot-specific policy, but these recorded runs do not implement it.
All 15 successful knottedness queries in the audit used this probe.
Their exact algebraic model expressions were retained and checked against
the assertions by Z3 itself.  This remains a trusted-solver computation,
not an independently checked proof certificate.  Neither a trace-zero
failure nor any timeout is accepted as an unknot decision.

The timing allowances are soft: the largest recorded audit call took
$375.926$ milliseconds with a requested $300$-millisecond allowance,
and the largest benchmark query took $288.326$ milliseconds with a
requested $250$-millisecond allowance.  The benchmark's reported query
time includes formula parsing and solver checks, but excludes the Z3
import and subsequent model extraction and validation.  Its compiler
times exclude construction of the input diagram.  The audit's per-call
times cover compilation and the solver wrapper.

The identical repeated compiler controls also show appreciable timing
variation: across 63 pairs, the median first-run/control ratio was 1.234,
with range 0.609--5.565.  Since the control always runs second, this
includes order and warm-up effects.  With three repetitions, fixed case
order, and a small convenience sample, the data justify neither a
statistical speedup claim nor an extrapolated asymptotic exponent.
They document correct completed queries and concrete representation
tradeoffs.  Renegar's theoretical complexity and this optional NLSAT
implementation remain separate statements.
